% GENERATED FILE. Edit canonical lessons, then run npm run generate:books.
% canonical-source-hash: fnv1a64:082db5967c7c2705
% canonical-lessons: SA-C138-vikrinati, SA-C138-karsati, SA-C138-nudati, SA-C138-bhramati, SA-C138-vardhate, SA-R138-first-pass-describing-words, SA-R138-second-pass-verbs

\chapter{Sells, Pulls, Pushes, Wanders, Grows}
\label{ch:sa-sells-pulls-pushes-wanders-grows}
\hlchaptermodality{138}

\begin{chapteropening}
\textbf{By the end of this chapter:} \emph{I can say sells, pulls, pushes, wanders and grows.}

Complete the last lesson of chapter 138: i can say sells, pulls, pushes, wanders and grows.
\end{chapteropening}

\section[vikrīṇāti]{\sk{विक्रीणाति} (vikrīṇāti) — sells}

\label{lesson:SA-C138-vikrinati}



\begin{quote}
Before the new one: say the Sanskrit for protects, then the Sanskrit for burns.
\end{quote}

\subsection*{You'll want to know: \sk{विक्रीणाति}}
\textbf{\sk{विक्रीणाति}} — \emph{vikrīṇāti} — \textquotedblleft{}sells\textquotedblright{}.

\begin{grammarlens}[title={What you've built}]
One more word for this chapter.
\end{grammarlens}

\subsection*{Your turn}
\begin{itemize}
\raggedright
  \item \emph{Say it:} \emph{vikrīṇāti}
  \item \emph{Say it:} \emph{vikrīṇāti}, once more
  \item \emph{Say it:} say \emph{vikrīṇāti}
  \item \emph{Recall:} say the Sanskrit for protects, then the Sanskrit for burns, then say \emph{vikrīṇāti} again
\end{itemize}

\begin{tcolorbox}[breakable,colback=teal!4,colframe=teal!35!black,title={Before you move on}]
What does \sk{विक्रीणाति} mean? (\textquotedblleft{}Sells\textquotedblright{}.) Say it once more.
\end{tcolorbox}
\section[karṣati]{\sk{कर्षति} (karṣati) — pulls, ploughs}

\label{lesson:SA-C138-karsati}



\begin{quote}
Before the new one: say the Sanskrit for burns, then the Sanskrit for sells.
\end{quote}

\subsection*{You'll want to know: \sk{कर्षति}}
\textbf{\sk{कर्षति}} — \emph{karṣati} — \textquotedblleft{}pulls, ploughs\textquotedblright{}.

\begin{grammarlens}[title={What you've built}]
One more word for this chapter.
\end{grammarlens}

\subsection*{Your turn}
\begin{itemize}
\raggedright
  \item \emph{Say it:} \emph{karṣati}
  \item \emph{Say it:} \emph{karṣati}, once more
  \item \emph{Say it:} say \emph{karṣati}
  \item \emph{Recall:} say the Sanskrit for burns, then the Sanskrit for sells, then say \emph{karṣati} again
\end{itemize}

\begin{tcolorbox}[breakable,colback=teal!4,colframe=teal!35!black,title={Before you move on}]
What does \sk{कर्षति} mean? (\textquotedblleft{}Pulls, ploughs\textquotedblright{}.) Say it once more.
\end{tcolorbox}
\section[nudati]{\sk{नुदति} (nudati) — pushes}

\label{lesson:SA-C138-nudati}



\begin{quote}
Before the new one: say the Sanskrit for sells, then the Sanskrit for pulls.
\end{quote}

\subsection*{You'll want to know: \sk{नुदति}}
\textbf{\sk{नुदति}} — \emph{nudati} — \textquotedblleft{}pushes\textquotedblright{}.

\begin{grammarlens}[title={What you've built}]
One more word for this chapter.
\end{grammarlens}

\subsection*{Your turn}
\begin{itemize}
\raggedright
  \item \emph{Say it:} \emph{nudati}
  \item \emph{Say it:} \emph{nudati}, once more
  \item \emph{Say it:} say \emph{nudati}
  \item \emph{Recall:} say the Sanskrit for sells, then the Sanskrit for pulls, then say \emph{nudati} again
\end{itemize}

\begin{tcolorbox}[breakable,colback=teal!4,colframe=teal!35!black,title={Before you move on}]
What does \sk{नुदति} mean? (\textquotedblleft{}Pushes\textquotedblright{}.) Say it once more.
\end{tcolorbox}
\section[bhramati]{\sk{भ्रमति} (bhramati) — wanders}

\label{lesson:SA-C138-bhramati}



\begin{quote}
Before the new one: say the Sanskrit for pulls, then the Sanskrit for pushes.
\end{quote}

\subsection*{You'll want to know: \sk{भ्रमति}}
\textbf{\sk{भ्रमति}} — \emph{bhramati} — \textquotedblleft{}wanders\textquotedblright{}.

\begin{grammarlens}[title={What you've built}]
One more word for this chapter.
\end{grammarlens}

\subsection*{Your turn}
\begin{itemize}
\raggedright
  \item \emph{Say it:} \emph{bhramati}
  \item \emph{Say it:} \emph{bhramati}, once more
  \item \emph{Say it:} say \emph{bhramati}
  \item \emph{Recall:} say the Sanskrit for pulls, then the Sanskrit for pushes, then say \emph{bhramati} again
\end{itemize}

\begin{tcolorbox}[breakable,colback=teal!4,colframe=teal!35!black,title={Before you move on}]
What does \sk{भ्रमति} mean? (\textquotedblleft{}Wanders\textquotedblright{}.) Say it once more.
\end{tcolorbox}
\section[vardhate]{\sk{वर्धते} (vardhate) — grows}

\label{lesson:SA-C138-vardhate}



\begin{quote}
Before the new one: say the Sanskrit for pushes, then the Sanskrit for wanders.
\end{quote}

\subsection*{You'll want to know: \sk{वर्धते}}
\textbf{\sk{वर्धते}} — \emph{vardhate} — \textquotedblleft{}grows\textquotedblright{}.

\begin{grammarlens}[title={What you've built}]
One more word for this chapter.
\end{grammarlens}

\subsection*{Your turn}
\begin{itemize}
\raggedright
  \item \emph{Say it:} \emph{vardhate}
  \item \emph{Say it:} \emph{vardhate}, once more
  \item \emph{Say it:} say \emph{vardhate}
  \item \emph{Recall:} say the Sanskrit for pushes, then the Sanskrit for wanders, then say \emph{vardhate} again
\end{itemize}

\begin{tcolorbox}[breakable,colback=teal!4,colframe=teal!35!black,title={Before you move on}]
What does \sk{वर्धते} mean? (\textquotedblleft{}Grows\textquotedblright{}.) Say it once more.
\end{tcolorbox}
\section[(dialogue)]{Describing words — a first pass}

\label{lesson:SA-R138-first-pass-describing-words}



\begin{quote}
Say the words for wanders and grows.
\end{quote}

\subsection*{Your turn}
Say these aloud:

\begin{itemize}
\raggedright
  \item foolish, lazy, strong, weak, clear
  \item black, grey, pink, clear, thick
  \item thin, crooked, straight, full, empty
  \item low, high, flies, swims, jumps
\end{itemize}

\begin{tcolorbox}[breakable,colback=teal!4,colframe=teal!35!black,title={Before you move on}]
Foolish? (\textbf{\sk{मूर्खः}}, \emph{mūrkhaḥ}.) Lazy? (\textbf{\sk{अलसः}}, \emph{alasaḥ}.) Strong? (\textbf{\sk{सबलः}}, \emph{sabalaḥ}.) Weak? (\textbf{\sk{दुर्बलः}}, \emph{durbalaḥ}.) Clear? (\textbf{\sk{स्वच्छः}}, \emph{svacchaḥ}.) Black? (\textbf{\sk{कृष्णः}}, \emph{kṛṣṇaḥ}.) Grey? (\textbf{\sk{धूसरः}}, \emph{dhūsaraḥ}.) Pink? (\textbf{\sk{पाटलः}}, \emph{pāṭalaḥ}.) Clear? (\textbf{\sk{स्पष्टः}}, \emph{spaṣṭaḥ}.) Thick? (\textbf{\sk{स्थूलः}}, \emph{sthūlaḥ}.) Thin? (\textbf{\sk{कृशः}}, \emph{kṛśaḥ}.) Crooked? (\textbf{\sk{वक्रः}}, \emph{vakraḥ}.) Straight? (\textbf{\sk{ऋजुः}}, \emph{ṛjuḥ}.) Full? (\textbf{\sk{पूर्णः}}, \emph{pūrṇaḥ}.) Empty? (\textbf{\sk{रिक्तः}}, \emph{riktaḥ}.) Low? (\textbf{\sk{नीचः}}, \emph{nīcaḥ}.) High? (\textbf{\sk{उन्नतः}}, \emph{unnataḥ}.) Flies? (\textbf{\sk{उड्डयते}}, \emph{uḍḍayate}.) Swims? (\textbf{\sk{तरति}}, \emph{tarati}.) Jumps? (\textbf{\sk{कूर्दति}}, \emph{kūrdati}.)
\end{tcolorbox}
\section[(dialogue)]{Verbs — a second pass}

\label{lesson:SA-R138-second-pass-verbs}



\begin{quote}
Say the words for cuts and breaks.
\end{quote}

\subsection*{Your turn}
Say these aloud:

\begin{itemize}
\raggedright
  \item cuts, breaks, places, sends, searches
  \item learns, teaches, forgets, climbs, climbs down
  \item reaches, wears, counts, wipes, ties
  \item sews, wins, hides, protects, burns
  \item sells, pulls, pushes, wanders, grows
\end{itemize}

\begin{tcolorbox}[breakable,colback=teal!4,colframe=teal!35!black,title={Before you move on}]
Cuts? (\textbf{\sk{छिनत्ति}}, \emph{chinatti}.) Breaks? (\textbf{\sk{भिनत्ति}}, \emph{bhinatti}.) Places? (\textbf{\sk{स्थापयति}}, \emph{sthāpayati}.) Sends? (\textbf{\sk{प्रेषयति}}, \emph{preṣayati}.) Searches? (\textbf{\sk{अन्विष्यति}}, \emph{anviṣyati}.) Learns? (\textbf{\sk{शिक्षते}}, \emph{śikṣate}.) Teaches? (\textbf{\sk{पाठयति}}, \emph{pāṭhayati}.) Forgets? (\textbf{\sk{विस्मरति}}, \emph{vismarati}.) Climbs? (\textbf{\sk{आरोहति}}, \emph{ārohati}.) Climbs down? (\textbf{\sk{अवरोहति}}, \emph{avarohati}.) Reaches? (\textbf{\sk{प्राप्नोति}}, \emph{prāpnoti}.) Wears? (\textbf{\sk{धारयति}}, \emph{dhārayati}.) Counts? (\textbf{\sk{गणयति}}, \emph{gaṇayati}.) Wipes? (\textbf{\sk{मार्जयति}}, \emph{mārjayati}.) Ties? (\textbf{\sk{बध्नाति}}, \emph{badhnāti}.) Sews? (\textbf{\sk{सीव्यति}}, \emph{sīvyati}.) Wins? (\textbf{\sk{जयति}}, \emph{jayati}.) Hides? (\textbf{\sk{गूहति}}, \emph{gūhati}.) Protects? (\textbf{\sk{रक्षति}}, \emph{rakṣati}.) Burns? (\textbf{\sk{ज्वलति}}, \emph{jvalati}.) Sells? (\textbf{\sk{विक्रीणाति}}, \emph{vikrīṇāti}.) Pulls? (\textbf{\sk{कर्षति}}, \emph{karṣati}.) Pushes? (\textbf{\sk{नुदति}}, \emph{nudati}.) Wanders? (\textbf{\sk{भ्रमति}}, \emph{bhramati}.) Grows? (\textbf{\sk{वर्धते}}, \emph{vardhate}.)
\end{tcolorbox}
